\section*{Bab \ref{ch:b1:optical-mineralogy} --- Mineralogi optik}

\begin{solution}{exo:b1:optical-mineralogy:1}
Kuarsa: $\qty{30000}{nm}\times0.009 = \qty{270}{nm}$, putih \omterm{def:b1:optical-mineralogy:interference-colour}{orde} pertama.
Kalsit: $\qty{30000}{nm}\times0.172 = \qty{5160}{nm}$, jauh di luar bagan:
putih \omterm{def:b1:optical-mineralogy:interference-colour}{orde} tinggi yang pucat, sedikit merah muda atau hijau.
\end{solution}

\begin{solution}{exo:b1:optical-mineralogy:2}
$I/I_{\max} = \sin^2 2\theta$: $0$ pada \ang{0}; $\sin^2\ang{45} = 0.5$ pada
\ang{22.5}; $1$ pada \ang{45}; $0.5$ pada \ang{67.5}; $0$ pada \ang{90}.
\end{solution}

\begin{solution}{exo:b1:optical-mineralogy:3}
\omterm{def:b1:rock-forming-minerals:groups}{Garnet} bersifat kubik, sehingga isotropik: cahaya meninggalkannya bergetar seperti ketika masuk, sepanjang polarisator, dan analisator menghentikannya; $\Delta = 0$ pada setiap posisi. Butir kuarsa yang dipotong tegak lurus \omterm{def:b1:optical-mineralogy:optic-axis}{sumbu optiknya} juga memiliki $n_1 = n_2$ untuk cahaya yang merambat sepanjang sumbu itu, sehingga $\Delta = 0$: ia pun tetap gelap. Dalam \omterm{def:b1:field-lab-methods-1:microscope}{cahaya terpolarisasi bidang} \omterm{def:b1:rock-forming-minerals:groups}{garnet} memiliki relief sangat tinggi ($1.83$ terhadap $1.54$), kuarsa hampir tanpa relief.
\end{solution}

\begin{solution}{exo:b1:optical-mineralogy:4}
$|n - 1.54|$: almandin $0.29$, forsterit $0.11$, ortoklas $0.02$, kuarsa $0.01$; relief dalam urutan itu, dari yang tertinggi ke yang terendah. Ortoklas memiliki indeks yang lebih rendah: relief negatif, dan \omterm{def:b1:optical-mineralogy:relief}{garis Becke-nya} bergerak keluar dari butir, ke dalam pemasang, ketika fokus dinaikkan.
\end{solution}

\begin{solution}{exo:b1:optical-mineralogy:5}
$\delta = 650/30000 = 0.022$. \omterm{def:b1:rock-forming-minerals:mineral}{Mineral} hijau \omterm{def:b1:optical-mineralogy:pleochroism}{pleokroik} dengan belahan pada \ang{56} dan \ang{124} dan \omterm{def:b1:optical-mineralogy:birefringence}{birefringensi} sekitar $0.020$: hornblende.
\end{solution}

\begin{solution}{exo:b1:optical-mineralogy:6}
$\tan\ang{30} = 0.5774$. Untuk kuarsa: $(0.013/1.54)\times0.5774 = \qty{0.00487}{rad}
= \ang{0.28}$. Untuk forsterit: $(0.111/1.54)\times0.5774 = \qty{0.0416}{rad} =
\ang{2.38}$. Perbandingan $0.111/0.013 = 8.5$. Secara tepat: $\sin r = 1.54\sin\ang{30}/1.651
= 0.4664$, $r = \ang{27.80}$, deviasi \ang{2.20}: rumus \omterm{def:b1:optical-mineralogy:interference-colour}{orde} pertama melebih-lebihkan sebesar \qty{8}{\percent} untuk selisih indeks yang sebesar ini.
\end{solution}

\begin{solution}{exo:b1:optical-mineralogy:7}
$270 + 540 = \qty{810}{nm}$, \omterm{def:b1:optical-mineralogy:interference-colour}{orde} kedua hijau, bila arah lambat butir dan keping sejajar (tundanya bertambah); $|270 - 540| =
\qty{270}{nm}$, putih \omterm{def:b1:optical-mineralogy:interference-colour}{orde} pertama lagi, bila arah lambat butir sejajar arah cepat keping.
\end{solution}

\begin{solution}{exo:b1:optical-mineralogy:8}
Pada tebal penuh $\Delta = \qty{30000}{nm}\times0.052 = \qty{1560}{nm}$. Dari tepi ($\Delta = 0$) ke dalam, \omterm{def:b1:optical-mineralogy:retardation}{retardasi} melewati 542 dan
\qty{1074}{nm} tetapi tidak \qty{1605}{nm}: dua pita merah-ungu.
\end{solution}

\begin{solution}{exo:b1:optical-mineralogy:9}
\omterm{def:b1:optical-mineralogy:extinction}{Sudut pemadaman} yang diukur pada sebuah butir bergantung pada bagaimana butir itu dipotong: hanya penampang dengan orientasi tertentu yang menunjukkan sudut terbesar, dan yang lain menunjukkan lebih sedikit, sampai nol. Nilai terbesar, sekitar \ang{37}, adalah yang dibandingkan: ia cocok dengan diopsid (sekitar \ang{38}), bukan hornblende (paling tinggi \ang{19}).
\end{solution}

\begin{solution}{exo:b1:optical-mineralogy:10}
Dengan analisator sepanjang $x$, komponen $x$ dari kedua gelombang adalah
$E_0\cos\theta\cdot\cos\theta$ dan $-E_0\sin\theta\,e^{-i\varphi}\cdot(-\sin\theta)$,
sehingga $A = E_0(\cos^2\theta + \sin^2\theta\,e^{-i\varphi})$ dan
\[
  |A|^2 = E_0^2\bigl(\cos^4\theta + \sin^4\theta + 2\sin^2\theta\cos^2\theta\cos\varphi\bigr)
  = E_0^2\bigl(1 - 2\sin^2\theta\cos^2\theta\,(1 - \cos\varphi)\bigr),
\]
menggunakan $\cos^4\theta + \sin^4\theta = 1 - 2\sin^2\theta\cos^2\theta$. Karena
$2\sin^2\theta\cos^2\theta = \frac12\sin^2 2\theta$ dan $1 - \cos\varphi =
2\sin^2(\varphi/2)$, $I_\parallel = I_0(1 - \sin^2 2\theta\sin^2(\pi\Delta/\lambda))$.
Pada setiap panjang gelombang, $I_\parallel + I_\perp = I_0$: warna di antara polarisator sejajar adalah komplemen dari warna di antara polarisator silang. Butir \qty{550}{nm}, magenta di antara polarisator silang, tampak hijau kekuningan pucat di antara polarisator sejajar.
\end{solution}

\begin{solution}{exo:b1:optical-mineralogy:11}
(a) $\Delta/\lambda$ berjalan dari $5000/700 = 7.14$ hingga $5000/400 = 12.5$; ia melewati bilangan bulat 8 sampai 12: lima nol. (b) Nol-nolnya berada pada $\lambda_k = \Delta/k$; dua yang berurutan berjarak $\Delta/k - \Delta/(k + 1) =
\Delta/(k(k + 1)) \approx \lambda^2/\Delta$, yaitu $550^2/5000 =
\qty{60.5}{nm}$ dekat \qty{550}{nm}. (c) Lima pita gelap sempit yang tersebar di spektrum menghilangkan sedikit dari setiap warna, dan sisanya diteruskan dengan bobot rata-rata setengah: mata melihat rona pucat yang hampir netral.
\end{solution}

\begin{solution}{exo:b1:optical-mineralogy:12}
Rata-rata $\sin^2 2\theta$ atas satu putaran adalah $\frac1{2\pi}\int_0^{2\pi}\sin^2 2\theta\,
\mathrm d\theta = \frac12$, sehingga intensitas rata-ratanya adalah
$\frac12 I_0\sin^2(\pi\Delta/\lambda)$. Dia menghitung $\delta = \Delta/30000$
bukan $\Delta/25000$ (dalam nm): semua nilainya adalah $25/30 = 0.83$ dari yang benar, terlalu rendah sebesar \qty{17}{\percent}.
\end{solution}

\begin{solution}{pb:b1:optical-mineralogy:1}
\textbf{1.} B: tepi gelap menandakan relief tinggi; batas A yang nyaris tak terlihat menandakan relief rendah.

\textbf{2.} A: $(1.544 + 1.553)/2 = 1.5485$, $|n - n_m| = 0.0085$. B:
$(1.664 + 1.672 + 1.694)/3 = 1.6767$, $|n - n_m| = 0.1367$.

\textbf{3.} Deviasi sebanding dengan $|n - n_m|$: $0.1367/0.0085 =
16.1$.

\textbf{4.} B memiliki indeks yang lebih tinggi daripada pemasang: relief positif.

\textbf{5.} \omterm{def:b1:optical-mineralogy:pleochroism}{Pleokroisme}; biotit.

\textbf{6.} Dua belahan pada \ang{56} dan \ang{124}: sebuah \omterm{def:b1:rock-forming-minerals:groups}{amfibol}.

\textbf{7.} $\qty{30000}{nm}\times0.009 = \qty{270}{nm}$: putih \omterm{def:b1:optical-mineralogy:interference-colour}{orde} pertama.

\textbf{8.} Setiap butir dipotong dengan orientasinya sendiri dan menunjukkan nilai $|n_1 - n_2|$ antara 0 dan $\delta$; hanya butir yang dipotong sejajar \omterm{def:b1:optical-mineralogy:optic-axis}{sumbu optik} yang menunjukkan $\delta$ itu sendiri. Warna tertinggi di antara banyak butir adalah yang dipakai.

\textbf{9.} $d = 342/0.009 = \qty{38000}{nm} = \qty{38}{\micro m}$.

\textbf{10.} $(38 - 30)/30 = \qty{26.7}{\percent}$.

\textbf{11.} $\sin^2(\pi\times342/550) = 0.861$; pada \qty{30}{\micro m},
$\sin^2(\pi\times270/550) = 0.999$.

\textbf{12.} $\qty{38000}{nm}\times0.008 = \qty{304}{nm}$ hingga
$\qty{38000}{nm}\times0.013 = \qty{494}{nm}$: putih hingga jingga-merah, semuanya di bawah akhir \omterm{def:b1:optical-mineralogy:interference-colour}{orde} pertama pada \qty{542}{nm}.

\textbf{13.} Antara 1074 dan \qty{1605}{nm}: \omterm{def:b1:optical-mineralogy:interference-colour}{orde} ketiga.

\textbf{14.} $\delta = 1140/38000 = 0.030$.

\textbf{15.} $1140/30000 = 0.038$, yang akan menyarankan \omterm{def:b1:rock-forming-minerals:groups}{olivin}. \omterm{def:b1:rock-forming-minerals:groups}{Olivin} kaya magnesium dan kuarsa bereaksi menghasilkan \omterm{def:b1:rock-forming-minerals:groups}{piroksen},
\ce{Mg2SiO4 + SiO2 -> 2 MgSiO3}; batuan yang telah mengkristal dengan kuarsa tidak seharusnya memiliki keduanya.

\textbf{16.} $0.030$ adalah \omterm{def:b1:optical-mineralogy:birefringence}{birefringensi} diopsid; relief tingginya, ketiadaan warnanya, dan dua belahannya di dekat \ang{87} cocok: B adalah klinopiroksen, diopsid.

\textbf{17.} $\qty{30000}{nm}\times0.030 = \qty{900}{nm}$: kuning \omterm{def:b1:optical-mineralogy:interference-colour}{orde} kedua.

\textbf{18.} Dengan menipiskan hingga \qty{8}{\micro m}; kuarsa kemudian paling tinggi menunjukkan \qty{270}{nm}, putih \omterm{def:b1:optical-mineralogy:interference-colour}{orde} pertama.

\textbf{19.} \omterm{def:b1:optical-mineralogy:twin}{Kembaran polisintetik}: D adalah \omterm{def:b1:rock-forming-minerals:groups}{plagioklas}.

\textbf{20.} Komposisi \omterm{def:b1:rock-forming-minerals:groups}{plagioklas} (kandungan anortitnya).

\textbf{21.} Sebuah \omterm{def:b1:rock-forming-minerals:groups}{amfibol} dengan \omterm{def:b1:optical-mineralogy:extinction}{sudut pemadaman} \ang{17}: hornblende
(\ang{15}--\ang{19}), bukan \omterm{def:b1:rock-forming-minerals:groups}{piroksen} (\ang{38} untuk diopsid).

\textbf{22.} $\qty{38000}{nm}\times0.020 = \qty{760}{nm}$: hijau \omterm{def:b1:optical-mineralogy:interference-colour}{orde} kedua.

\textbf{23.} Ia dipotong tegak lurus \omterm{def:b1:optical-mineralogy:optic-axis}{sumbu optik}: kedua indeksnya sama untuk cahaya yang merambat sepanjang sumbu itu, $\Delta = 0$, dan ia tetap gelap seperti butir isotropik.

\textbf{24.} A kuarsa, B diopsid, C biotit, D \omterm{def:b1:rock-forming-minerals:groups}{plagioklas}, E hornblende: batuan beku menengah dengan sedikit kuarsa, seperti diorit kuarsa.

\textbf{25.} Sayatan itu setebal \textbf{\qty[text-series-to-math=true]{38}{\micro m}}.
\end{solution}
